\documentclass[10pt,a4paper]{amsart}
\usepackage[margin=19mm]{geometry}
\usepackage{amsmath,amssymb,mathtools}
\usepackage[scr=rsfs,cal=euler]{mathalfa}
\usepackage{xfrac}
\usepackage[unicode,hidelinks,bookmarks=false]{hyperref}
\newcommand{\ie}{i.e.\ }
\newcommand{\cf}{cf.\ }
\newcommand{\resp}{respectively\ }
\newcommand{\Subsection}{\subsection}
\providecommand{\nobreakdash}{\nobreak\hbox{-}}
\setlength{\emergencystretch}{2em}
\multlinegap0pt
\usepackage{mathtools}
\usepackage{amssymb}
\usepackage[shortlabels]{enumitem}
\usepackage{dsfont}
\usepackage{stackengine}
\usepackage{tikz}
\newtheorem{theorem}{Theorem}
\newtheorem{lemma}{Lemma}
\DeclareMathOperator{\RCvtex}{RC}
\title{On Nondefinability of Interior-Connectedness\break via the Contact Relation}
\author{Rafa{\l} Gruszczy\'nski, Paula Mench\'on}
\date{Source: arXiv:2503.08691v2 (CC BY 4.0)}
\begin{document}
\maketitle

\noindent\emph{Source and licence.} On Nondefinability of Interior-Connectedness\break via the Contact Relation. Version arXiv:2503.08691v2 (CC BY 4.0), distributed by arXiv under the Creative Commons Attribution 4.0 International licence, http://creativecommons.org/licenses/by/4.0/, as recorded in the arXiv licence metadata for this exact version and shown on the abstract page https://arxiv.org/abs/2503.08691v2. Full licence notice: \texttt{licenses/arxiv-2503.08691-CC-BY-4.0.txt}.\par\smallskip

\noindent\textit{Editorial scope.} Counted unit: the article's lemma lem:4-points (source0.tex lines 866--872). The article's complete original proof of it is delivered with it (source0.tex lines 874--910, one \texttt{proof} environment of 37 lines). No mathematics is written, repaired or re-derived: the statement and the proof are the article's own lines, in the article's own notation, and no retained line is substituted or re-flowed.  Retained uncounted as the source-local context the proof needs: lines 764--775; lines 845--857.  Nothing was omitted from any retained span, so the package carries no omission record.  No retained line contains a citation, so the wrapper carries no bibliography and no bibliography entry.  No retained line contains a figure environment, an image-inclusion command or drawing code, so no image is delivered and the package declares no asset dependency.  Editorial formatting only, in addition: an AMS \texttt{amsart} preamble that restates the article's own declarations and macros used by the retained lines, namely the environments \texttt{theorem}, \texttt{lemma} on separate counters, together with the standard packages those lines need.  The retained environments print in this document as Theorem 1, Lemma 1 in order of appearance, whereas the article prints the same items with the numbering of its own full document.  The complete native source of the article as distributed in the arXiv e-print for this exact version is bundled unchanged as \texttt{sources/174599/source0.tex}.
\begin{theorem}%
%%LEAP%%%\label{thm3.1}
\label{th:8-elements}
Let $  \langle X_{1},\tau _{1}  \rangle $ and
$  \langle X_{2},\tau _{2}  \rangle $ be topological spaces. If
$f$ is a BCA-isomorphism of the contact algebras
$\mathfrak{B}_1\coloneqq  \langle \RCvtex (\tau _{1}),\mathord{\mathrel{\mathsf{C}}_{\tau _{1}}}
  \rangle $ and
$\mathfrak{B}_2\coloneqq  \langle \RCvtex (\tau _{2}),\mathord{\mathrel{\mathsf{C}}_{\tau _{2}}}
  \rangle $, but $f$ does not preserve interior-connectedness, then
both algebras must have at least eight elements.
\end{theorem}

\begin{theorem}%
%%LEAP%%%\label{thm4.1}
\label{th:minimality-theorem}
Let $  \langle X_{1},\tau _{1}  \rangle $ and
$  \langle X_{2},\tau _{2}  \rangle $ be topological spaces. If
$f$ is a BCA-isomorphism of the algebras
$\mathfrak{B}_{1}\coloneqq   \langle \RCvtex (\tau _{1}),\mathord{
\mathrel{\mathsf{C}}_{\tau _{1}}}  \rangle $ and
$\mathfrak{B}_{2}\coloneqq   \langle \RCvtex (\tau _{2}),\mathord{
\mathrel{\mathsf{C}}_{\tau _{2}}}  \rangle $, but $f$ does not preserve
interior-connectedness, then the minimal number of points of the underlying
spaces is five for one of the spaces and four for the other.
\end{theorem}

\begin{lemma}%
%%LEAP%%%\label{lem4.2}
\label{lem:4-points}
If $\mathfrak{B}_{1}$ and $\mathfrak{B}_{2}$ and $f$ are like in Theorem~\ref{th:minimality-theorem},
then each of the underlying topological spaces must have at least four
points.
\end{lemma}

\begin{proof}
Assume we have isomorphic BCAs $\mathfrak{B}_{1}$ and
$\mathfrak{B}_{2}$ that have different properties of interior-connectedness.
Let $f\colon \mathfrak{B}_{1}\to \mathfrak{B}_{2}$ be their BCA-isomorphism.
Suppose the first is the algebra of regular closed sets of a space
$\langle X_{1},\tau _{1}\rangle $ whose domain has precisely three points:
$\mathfrak{B}_{1}\coloneqq   \langle \RCvtex (\tau _{1}),\mathrel{
\mathsf{C}}_{1}  \rangle $. In light of Theorem~\ref{th:8-elements},
it must be the case that $\RCvtex (\tau _{1})=2^{X_{1}}$, so all subsets of
$X_{1}$ are closed, and $X_{1}$ is a discrete space. This also means that
$\RCvtex (\tau _{1})$ has the minimal contact relation. In this case, the only
nonempty interior-connected subsets of $X_{1}$ are the singletons.

Suppose $\langle X_{2},\tau _{2}\rangle $ is a topological space whose
regular closed algebra
$\mathfrak{B}_{2}\coloneqq   \langle \RCvtex (\tau _{2}),\mathord{
\mathrel{\mathsf{C}}_{2}}  \rangle $ is isomorphic to
$\mathfrak{B}_1$, that is, $\RCvtex (\tau _{2})$ has exactly eight elements.
Since $f$ is an isomorphism of contact algebras, we have that
%
\begin{equation*}
A\cap D\neq \emptyset \qquad \text{iff}\qquad f(A)\mathrel{
\mathsf{C}}_{2} f(D).
\end{equation*}
%
Let $\cdot $ be the operation of meet of $\RCvtex (\tau _{2})$. Then we have
that $f(A)\mathrel{\mathsf{C}}_{2} f(D)$ entails that
$A\cap D\neq \emptyset $, so $f(A)\cdot f(D)\neq \emptyset $, as $f$ is
a Boolean homomorphism. This implies that the contact relation of
$\RCvtex (\tau _{2})$ must also be minimal. But this means that
$\RCvtex (\tau _{2})$ coincides with the algebra of clopen subsets of the space
$X_{2}$. So, its atoms are the only connected components of the space that
are also interior-connected. Therefore, $\mathfrak{B}_{1}$ and
$\mathfrak{B}_{2}$ have the same notions of interior-connectedness, which
is a contradiction. So, both spaces must have at least four points each.
This ends the proof.
\end{proof}

\end{document}
